% CODEX: wstaw bezpośrednio przed \section{Metryki lorentzowskie: Schwarzschild, Kerr i FLRW}
% w 11-metryka-dlugosc-miara.tex. Nie powtarzaj definicji metryk z rozdziału 30b.
\begin{przyklad}[Objętość i anizotropia w geometriach Nil i Sol]
\label{prz:thur-dop-metryki-nil-sol}
W rozdziale o geometryzacji przedstawimy dwie jednorodne metryki
trójwymiarowe, które nie mają stałej krzywizny sekcyjnej:
$g_N=(\dd x)^2+(\dd y)^2+(\dd z-x\,\dd y)^2$ oraz
$g_S=e^{-2t}(\dd x)^2+e^{2t}(\dd y)^2+(\dd t)^2$.
Nie przepisujemy tu ich konstrukcji grupowych; zbadamy wielkości,
które zależą tylko od współczynników metryki.

Dla $g_N$ macierz w mapie $(x,y,z)$ ma postać
\[
G_N=\begin{pmatrix}1&0&0\\0&1+x^2&-x\\0&-x&1\end{pmatrix},
\qquad \det G_N=1.
\]
Zatem riemannowska forma objętości przy orientacji $(x,y,z)$ to
$\dd x\wedge\dd y\wedge\dd z$, mimo że współczynniki metryki
zależą od $x$. W modelu Sol macierz jest diagonalna,
$G_S=\operatorname{diag}(e^{-2t},e^{2t},1)$,
więc także $\det G_S=1$ i
$\mathrm{vol}_{g_S}=\dd x\wedge\dd y\wedge\dd t$.
Dla poziomego wektora $v=a\partial_x+b\partial_y$ w Sol
mamy $|v|_{g_S}^2=e^{-2t}a^2+e^{2t}b^2$:
odcinki równoległe do dwóch osi podlegają przeciwnym zmianom
skali, chociaż element objętości jest niezmienny.
W Nil z kolei forma $\theta=\dd z-x\,\dd y$ mierzy składową
wektora w wyróżnionym kierunku pionowym.
Wzór na objętość nie rozstrzyga zatem ani o płaskości,
ani o izotropii metryki; dokładny rachunek krzywizn
przeprowadzimy w Rozdziale~\ref{ch:krzywizna}.
\end{przyklad}
